\documentclass[a4paper,12pt]{article}

\usepackage{a4wide}
\usepackage{amsfonts}
\usepackage{amsmath}
\usepackage{amssymb}
\usepackage{lipsum}

\usepackage{graphicx}  % For including images

\usepackage{xcolor}  % For a colorfull presentation
\usepackage{listings}  % For presenting code

\usepackage{hyperref}

% Definition of a style for code, matter of taste
\lstdefinestyle{mystyle}{
  language=Python,
  basicstyle=\ttfamily\footnotesize,
  backgroundcolor=\color[HTML]{F7F7F7},
  rulecolor=\color[HTML]{EEEEEE},
  identifierstyle=\color[HTML]{24292E},
  emphstyle=\color[HTML]{005CC5},
  keywordstyle=\color[HTML]{D73A49},
  commentstyle=\color[HTML]{6A737D},
  stringstyle=\color[HTML]{032F62},
  emph={@property,self,range,True,False},
  morekeywords={super,with,as,lambda},
  literate=%
    {+}{{{\color[HTML]{D73A49}+}}}1
    {-}{{{\color[HTML]{D73A49}-}}}1
    {*}{{{\color[HTML]{D73A49}*}}}1
    {/}{{{\color[HTML]{D73A49}/}}}1
    {=}{{{\color[HTML]{D73A49}=}}}1
    {/=}{{{\color[HTML]{D73A49}=}}}1,
  breakatwhitespace=false,
  breaklines=true,
  captionpos=b,
  keepspaces=true,
  numbers=none,
  showspaces=false,
  showstringspaces=false,
  showtabs=false,
  tabsize=4,
  frame=single,
}
\lstset{style=mystyle}

\begin{document}
\title{Machine Learning A (2026)\\Home Assignment 1}
\author{\color{red}Ludovico Maria Spitaleri, DMH249}
\date{}
\maketitle

% Please leave the table of contents as is, for the ease of navigation for TAs
\tableofcontents % Generates the table of contents
\newpage % Start a new page after the table of contents

\section{Make Your Own (10 points)}

% Enumerator environment, add items as needed
\begin{enumerate}
    \item The features I would include in the input space $\mathcal{X}$ would be:
        \begin{itemize}
            \item The current GPA of the student $x_\text{gpa} \in [0, 1]$, which provides a good estimate
            of the general performance of the candidate.
            \item The fields of study $x_\text{field} \in \{0, 1\}^k$, where $k$ is
            the number of fields of study, since students with knowledge in related
            fields are more likely to perform well in the course.
            \item $n$ pairs of features, one for each existing course, containing
            whether the student attended it $x_\text{attended} \in \{0, 1\}$ and the final grade he obtained,
            with a value of $x_\text{grade} \in [0, 1]$ so that it is independent of any specific grading scale.
            This could help detect more specific correlations between a certain final
            grade in the ML course and other subjects than just including the GPA,
            that averages all subjects.
        \end{itemize}
        Combining all that, the final input space would be
        $\mathcal{X} = [0, 1] \times \{0, 1\}^k \times \{0, 1\}^n \times [0, 1]^n$.
    \item The label space would be $\mathcal{Y} = [0, 1]$ such that
        the value can then easily be converted into any preferred grading scale.
    \item Considering the label space $\mathcal{Y} = [0, 1]$, this
        would give an infinite number of options, making this problem a regression one.
        A squared error loss function $\ell(y, \hat{y}) = (y - \hat{y})^2$
        makes the most sense since we can calculate the actual difference between
        the predicted and the real grade.
    \item The distance measure used for this problem could be a simple Euclidean
        distance, given by the norm of the difference between the two vectors
        representing the student's features $d(x, x') = \left|\left|x - x'\right|\right|$.
    \item The most important rule to evaluate a model is to avoid biased results
        caused by the reuse of data. So my evaluation process would consist of 3 phases:
        \begin{enumerate}
            \item Take the first part of my dataset $\mathcal{S}_1$ to train the model and get a
                hypothesis class $\mathcal H = \{ h_i: \mathcal{X} \rightarrow \mathcal{Y} \}$
                for $i = 1, \dots, M$.
            \item Take the second part of my dataset $\mathcal{S}_2(|\mathcal{S}_2| = N)$
                to select the hypothesis with the lowest average error.
                $$
                \begin{aligned}
                r_i &= \mathcal{R}_{S_2}(h_i) = \frac{1}{N} \sum_{j=1}^{N} \ell(y_j, h_i(x_j)) & \forall\, i = 1, \dots, M \\
                h_{i^*} &\in \mathcal{H} \text{ with } i^* = \operatorname{arg\,min}_{i} r_i &
                \end{aligned}
                $$
            \item Take the third and final part of my dataset $\mathcal{S}_3$ to evaluate the selected
                hypothesis $h_{i^*}$ by seeing the average error on new and unbiased data.
        \end{enumerate}
    \item The main issue that I could expect is that the university could update,
        add or remove courses, making the model obsolete fast. A solution to that
        would be to train the model on a fixed schedule to be up to date with new
        data, but if a course was changed or added recently it would be difficult
        to have enough data to re-train the model.
\end{enumerate}

\section{Digits Classification with K Nearest Neighbors (40 points)}

\subsection{Task \#1}

\subsection{Task \#2 (optional, not for submission)}





\section{Regression (50 points)}

\subsection{Task 1 (5 points)}

\subsection{Task 2 (10 points)}

\subsection{Task 3 (10 points)}

\subsection{Task 4 (5 points)}

\subsection{Task 5 (10 points)}

\subsection{Task 6 (10 points)}




% Placeholder figure
% \begin{figure}[htbp]
%     \centering
%     \includegraphics[width=0.5\linewidth]{example-image}
%     \caption{\lipsum[1][1-3]} % Dummy caption generated using lipsum
%     \label{fig:placeholder}
% \end{figure}



% Placeholder code
\begin{lstlisting}
# Example code
# Creating an example array
data = np.array([5, 2, 8, 1, 6])

# Calculating cumulative sum using cumsum
cumulative_sum = np.cumsum(data)
\end{lstlisting}



\end{document}
